\documentclass[a4paper,10pt]{article}
\usepackage{lamia-pedagogy}

\begin{document}

% =========================================================================
% PAGE 1 : Fiche de Cours - La Relative par "qui"
% =========================================================================
\ficheHeader{1 Heure}{\makebox[2.5cm]{\dotfill}}{1\textsuperscript{ère} AM}{2026 / 2027}

\vspace{-3pt}
\competenceBanner{\faBookOpen\ Projet 01 / Séquence 01 — Séance 05 : Grammaire}

\vspace{-2pt}
\modaliteSupportBar{Observation dirigée puis exercices écrits}{Support : Phrases d'observation / Cahier de grammaire}

\vspace{1pt}

% Phrases supports
\begin{tcolorbox}[
  enhanced,
  colback=primarylight!40,
  colframe=primary,
  arc=2mm,boxrule=0.8pt,
  top=4pt,bottom=4pt,left=8pt,right=8pt,
  title={\small\textbf{\textcolor{white}{\faEye\ Phrases d'observation}}},
  coltitle=white,colbacktitle=primary
]
{\small
\begin{enumerate}[label=\textbf{\alph*)},leftmargin=*,noitemsep,topsep=1pt]
  \item La douche est une \textbf{activité} \textcolor{softred}{\textbf{qui}} permet d'éliminer la transpiration.
  \item Elle utilise un \textbf{savon} \textcolor{softred}{\textbf{qui}} contient des ingrédients hydratants pour la peau.
  \item Les dentifrices soignent les \textbf{dents} \textcolor{softred}{\textbf{qui}} sont attaquées par les bactéries.
\end{enumerate}
}
\end{tcolorbox}

\vspace{3pt}

% Analyse et Définition de la Règle
\begin{minipage}[t]{0.485\textwidth}
\begin{tcolorbox}[
  enhanced,
  colback=secondarylight!50,
  colframe=secondary,
  arc=2mm,boxrule=0.8pt,
  top=4pt,bottom=4pt,left=7pt,right=7pt,
  title={\small\textbf{\textcolor{secondary!90!black}{\faStar\ Le Rôle de \guillemotleft\ qui \guillemotright}}},
  colbacktitle=white,coltitle=secondary!90!black,
  attach boxed title to top left={yshift=-2mm,xshift=4mm},
  boxed title style={boxrule=0.6pt,colframe=secondary,arc=1.2mm,top=1pt,bottom=1pt}
]
\vspace{2pt}
{\small
\textbf{\guillemotleft\ Qui \guillemotright\ est un pronom relatif :}
\begin{itemize}[leftmargin=*,noitemsep,topsep=1pt]
  \item Il sert à \textbf{relier deux phrases simples} pour en faire une phrase complexe.
  \item Il permet d'\textbf{éviter la répétition} d'un mot.
  \item Le mot placé juste avant \guillemotleft\ qui \guillemotright\ s'appelle \textbf{l'antécédent}.
\end{itemize}
}
\end{tcolorbox}
\end{minipage}%
\hfill
\begin{minipage}[t]{0.485\textwidth}
\begin{tcolorbox}[
  enhanced,
  colback=accentlight!35,
  colframe=accent,
  arc=2mm,boxrule=0.8pt,
  top=4pt,bottom=4pt,left=7pt,right=7pt,
  title={\small\textbf{\textcolor{accent!90!black}{\faStar\ La Structure de la PSR}}},
  colbacktitle=white,coltitle=accent!90!black,
  attach boxed title to top left={yshift=-2mm,xshift=4mm},
  boxed title style={boxrule=0.6pt,colframe=accent,arc=1.2mm,top=1pt,bottom=1pt}
]
\vspace{2pt}
{\small
\textbf{Proposition Subordonnée Relative :}
\begin{itemize}[leftmargin=*,noitemsep,topsep=1pt]
  \item La proposition principale apporte l'information de base.
  \item La proposition introduite par \guillemotleft\ qui \guillemotright\ s'appelle la \textbf{P.S.R.} (Proposition Subordonnée Relative).
  \item Elle complète et précise l'antécédent.
\end{itemize}
}
\end{tcolorbox}
\end{minipage}

\vspace{3pt}

% Formule Visuelle Clé
\begin{tcolorbox}[
  enhanced,
  colback=white,
  colframe=cardborder,
  arc=2mm,boxrule=0.7pt,
  top=4pt,bottom=4pt,left=8pt,right=8pt,
  title={\small\textbf{\textcolor{primary}{\faPuzzlePiece\ Formule visuelle de la phrase complexe}}},
  coltitle=primary,colbacktitle=primarylight,titlerule=0.4pt
]
\centering
\vspace{1pt}
{\small
\fbox{\textbf{Proposition Principale}} \quad $+$\quad
\textcolor{secondary}{\fbox{\textbf{Antécédent}}} \quad $+$\quad
\textcolor{softred}{\fbox{\textbf{qui}}} \quad $+$\quad
\textcolor{accent!90!black}{\fbox{\textbf{Proposition Subordonnée Relative (P.S.R.)}}}
}
\vspace{2pt}
\par\noindent{\scriptsize\textit{Exemple : J'admire les \textbf{\textcolor{secondary}{médecins}} \textbf{\textcolor{softred}{qui}} \textbf{\textcolor{accent!90!black}{soignent les malades avec dévouement}}.}}
\end{tcolorbox}

\newpage

% =========================================================================
% PAGE 2 : Fiche d'Exercices d'Application
% =========================================================================
\ficheEleveHeader

\vspace{-3pt}
\competenceBanner{\faPenFancy\ Fiche d'Exercices : La Relative par \guillemotleft\ qui \guillemotright}

\vspace{2pt}

% ACTIVITÉ 1 : Repérage
\begin{tcolorbox}[
  enhanced,
  colback=white,
  colframe=cardborder,
  arc=2mm,boxrule=0.8pt,
  top=4pt,bottom=4pt,left=8pt,right=8pt,
  title={\small\textbf{\textcolor{primary}{\faHighlighter\ Activité 01 : Souligne la P.S.R., encadre le pronom relatif et entoure l'antécédent}} \hfill \badge[primary]{Analyse}},
  coltitle=primary,colbacktitle=primarylight,titlerule=0.5pt,titlerule style=cardborder
]
{\small
\begin{enumerate}[label=\textbf{\arabic*.},leftmargin=*,noitemsep,topsep=2pt]
  \item Le brossage est une habitude \textbf{qui} protège nos dents des caries.
  \item Les enfants écoutent les conseils du médecin \textbf{qui} explique l'importance de l'hygiène.
  \item Les microbes \textbf{qui} se trouvent sur nos mains sont invisibles à l'\oe il nu.
  \item Nous utilisons une serviette propre \textbf{qui} absorbe rapidement l'eau.
\end{enumerate}
}
\end{tcolorbox}

\vspace{2pt}

% ACTIVITÉ 2 : Évitement de répétition
\begin{tcolorbox}[
  enhanced,
  colback=white,
  colframe=cardborder,
  arc=2mm,boxrule=0.8pt,
  top=4pt,bottom=4pt,left=8pt,right=8pt,
  title={\small\textbf{\textcolor{primary}{\faLink\ Activité 02 : Évite la répétition en employant le pronom relatif \guillemotleft\ qui \guillemotright}} \hfill \badge[secondary]{Transformation}},
  coltitle=primary,colbacktitle=primarylight,titlerule=0.5pt,titlerule style=cardborder
]
{\small
\begin{enumerate}[label=\textbf{\arabic*.},leftmargin=*,noitemsep,topsep=2pt]
  \item Le docteur examine l'élève. L'élève a mal aux dents.\\
    $\longrightarrow$ \drule{\linewidth}
  \item Le savon élimine les bactéries. Les bactéries provoquent des infections.\\
    $\longrightarrow$ \drule{\linewidth}
  \item Maman achète un dentifrice. Ce dentifrice renforce les gencives.\\
    $\longrightarrow$ \drule{\linewidth}
\end{enumerate}
}
\end{tcolorbox}

\vspace{2pt}

% ACTIVITÉ 3 : Compléter une PSR
\begin{tcolorbox}[
  enhanced,
  colback=white,
  colframe=cardborder,
  arc=2mm,boxrule=0.8pt,
  top=4pt,bottom=4pt,left=8pt,right=8pt,
  title={\small\textbf{\textcolor{primary}{\faEdit\ Activité 03 : Complète les phrases avec la P.S.R. qui convient}} \hfill \badge[emerald]{Complétion}},
  coltitle=primary,colbacktitle=primarylight,titlerule=0.5pt,titlerule style=cardborder
]
{\small
\textit{Propositions : qui causent des caries / qui rafraîchit l'haleine / qui préserve la santé.}\\
\begin{enumerate}[label=\textbf{\alph*)},leftmargin=*,noitemsep,topsep=2pt]
  \item L'hygiène des mains est une règle \drule{6.5cm}.
  \item Les enfants évitent les bonbons sucrés \drule{6cm}.
  \item Le dentifrice à la menthe est une pâte \drule{6cm}.
\end{enumerate}
}
\end{tcolorbox}

\vspace{2pt}

% ACTIVITÉ 4 : Production
\begin{tcolorbox}[
  enhanced,
  colback=white,
  colframe=cardborder,
  arc=2mm,boxrule=0.8pt,
  top=4pt,bottom=4pt,left=8pt,right=8pt,
  title={\small\textbf{\textcolor{primary}{\faPenSquare\ Activité 04 : Production personnelle}} \hfill \badge[softred]{Production}},
  coltitle=primary,colbacktitle=primarylight,titlerule=0.5pt,titlerule style=cardborder
]
{\small
Rédige une phrase complexe avec le pronom relatif \textbf{\guillemotleft\ qui \guillemotright} pour expliquer pourquoi il faut se laver le corps régulièrement :\\
$\longrightarrow$ \drule{\linewidth}
}
\end{tcolorbox}

\end{document}
